\section{The lowest Weil eigenvector and the simplicity question}
\label{sec:weil-parity}
This section connects the lowest eigenvector in the programme of Connes,
Consani and Moscovici to the product over distinct zero ordinates, and then
to the multiplicities of the Riemann xi function.

Use the Fourier transform
$\widehat f(z)=\int_{\mathbb R}f(x)e^{-izx}\,dx$ and inner products linear
in their first argument. Write $\Gamma$ for the distinct real ordinates,
$m_\gamma$ for their multiplicities, and
\[
 q(f,g)=\sum_{\gamma\in\Gamma}
 m_\gamma\widehat f(\gamma)\overline{\widehat g(\gamma)}.
\]
The Riemann hypothesis, now provided by the input theorem, identifies
this expression with the positive Weil form. Its global completion is
$H=\ell^2(\Gamma,m)$ under sampling. The local completion $H_L$ of
$C_c^\infty(-L,L)$ embeds in $H$ and compactly in $L^2(-L,L)$.
The closed Weil operator $A_L$ has positive eigenvalues
$a_1(L)\le a_2(L)\le\cdots$, listed with multiplicity.
These local form and compact-resolvent facts are established in
\cite[\S3]{ccm2025}; the sampling realization and its unshifted inverse
compression are recorded in \cite[\S2.3 and Theorem 3.3]{suzuki}.
The form domain consists of zero-extended functions for which
$\int(1+\log(2+|t|))|\widehat f(t)|^2\,dt$ is finite. In this norm,
shrinking support by a dilation and then mollifying proves that compact
smooth functions are a form core. Dilation continuity follows first on
Schwartz functions, then by density and uniform boundedness for dilation
factors between $1/2$ and $1$.

Let $J_L:H_L\to L^2(-L,L)$ be inclusion, $\Pi_L:H\to H_L$ the
orthogonal projection, and $C$ multiplication by $\gamma^{-1}$ on $H$.
There is no zero at ordinate zero: the alternating-series identity gives
$(1-\sqrt2)\zeta(1/2)>0$. The order-one zero count gives
$\sum_\gamma m_\gamma/\gamma^2<\infty$, so $C$ is Hilbert--Schmidt.
Put $C_L=\Pi_LC\Pi_L$. Reflection decomposes $H_L$ into even and odd
parts, in which
\[
 C_L=\begin{pmatrix}0&X_L^*\\X_L&0\end{pmatrix},\qquad
 J_L^*J_L=\begin{pmatrix}T_{L,+}&0\\0&T_{L,-}\end{pmatrix}.
\]
Let $a_e(L)$ and $a_o(L)$ denote the minima in these two parts.
All the adjoints and projections in these formulas use $q$.

\begin{proposition}[The parity comparison]\label{prop:parity-comparison}
At a window with $a_e(L)<a_o(L)$ the lowest eigenvalue is simple and
even, and its eigenvector $e_L$ has nonzero integral. Put
\[
 \alpha_L=a_e(L),\quad M_{L,\pm}=\alpha_LT_{L,\pm},\quad
 S_{L,\pm}=I-M_{L,\pm},\quad
 \kappa_L=1-a_e(L)/a_o(L).
\]
On $H_L/\mathbb Ce_L$, with the form
$q-\alpha_L\langle\cdot,\cdot\rangle_2$, the compact primitive is
unitarily represented by the restriction off $e_L$ of
\[
 B_L=\begin{pmatrix}
 0&S_{L,+}^{1/2}X_L^*S_{L,-}^{-1/2}\\
 S_{L,-}^{-1/2}X_LS_{L,+}^{1/2}&0
 \end{pmatrix}.
\]
It is injective and selfadjoint. The reciprocals of its nonzero
eigenvalues are precisely the zeros of $\widehat e_L$, all of which
are real and simple. Moreover,
\begin{align}
 \|B_L-C_L\|_{\mathcal S_2}
 \le\sqrt2\bigl(&\kappa_L^{-1/2}
       \|M_{L,-}X_L\|_{\mathcal S_2}
       +\|X_LM_{L,+}\|_{\mathcal S_2}\bigr).
 \label{eq:parity-error}
\end{align}
Along every sequence $L\to\infty$ on which $\inf\kappa_L>0$,
\[
 \|B_L-C\|_{\mathcal S_2}\longrightarrow0,\qquad
 \frac{\widehat e_L(z)}{\widehat e_L(0)}\longrightarrow
 \Delta_\Gamma(z):=\prod_{\gamma>0,\ \gamma\in\Gamma}
              (1-z^2/\gamma^2)
\]
locally uniformly in $z$. Each distinct positive ordinate appears once
in this product.
\end{proposition}

\begin{proof}
Fix a window and write $\ell(f)=\int f$. Coercivity makes $\ell$ a
bounded functional on $H_L$. For a compact smooth mean-zero function,
\[
 C_Lf=\mathcal Pf,\qquad
 (\mathcal Pf)(x)=i\int_{-L}^x f(y)\,dy.
\]
Indeed its Fourier transform is $\widehat f(t)/t$, and its support
remains in the interval. This identity extends to every $f\in\ker\ell$:
approximate in the form norm by compact smooth functions, correct their
integrals with a fixed compact smooth function of integral one, and use
the boundedness of $C_L$. Their primitives also converge in
$H^1_0(-L,L)$, so the two limits agree in $L^2$.
Thus $DC_Lf=f$ on $\ker\ell$, where $D=-i\partial_x$, and
$C_LDg=g$ for compact smooth $g$.

An even bottom vector $u$ with integral zero would give a nonzero odd
primitive $v=C_Lu$. Its integral is zero by parity, so $w=C_Lv$ also
has compact support and $Dv=u$, $Dw=v$. The bottom eigenvector identity
and integration by parts give
\[
 q(v)=q(u,w)=\alpha_L\langle u,w\rangle_2
       =\alpha_L\langle Dv,w\rangle_2
       =\alpha_L\|v\|_2^2.
\]
This contradicts the strictly larger odd minimum. Two independent even
bottom vectors would have a nonzero mean-zero linear combination.
Consequently the bottom line is simple, even, and has nonzero integral.
Choose its $q$-unit vector $e$. On its orthogonal complement,
$q_\alpha=q-\alpha_L\langle\cdot,\cdot\rangle_2$ is an equivalent
Hilbert norm because $J_L^*J_L$ is compact and its largest eigenvalue
is simple.

Integration by parts for two mean-zero functions yields
$\langle C_Lf,g\rangle_2=\langle f,C_Lg\rangle_2$.
Hence $X_LM_{L,+}-M_{L,-}X_L$ vanishes on the even kernel of $\ell$.
Its value at $e$ is $S_{L,-}X_Le$, and therefore
\[
 X_LM_{L,+}-M_{L,-}X_L
   =S_{L,-}X_Le\,\frac{\ell}{\ell(e)}.
\]
Define $K_Lf=C_L(f-e\ell(f)/\ell(e))$. It kills $e$ and is the
compact primitive of the mean-zero expression in parentheses.
In quotient coordinates orthogonal to $e$, conjugation by
$S_{L,+}^{1/2}\oplus S_{L,-}^{1/2}$ gives the displayed $B_L$;
the commutator identity rewrites its lower-left block as
$S_{L,-}^{-1/2}X_LS_{L,+}^{1/2}$. The two blocks are adjoints.
No inverse of $S_{L,+}$ on its kernel is taken.

For injectivity, $C_LDg=g$ and the density of compact smooth $g$ first
show that $C_L$ is injective. If $K_Lf=ce$ and $c\ne0$, then
$e=C_Lg$ for a mean-zero $g$; injectivity and reflection make $g$ odd.
The primitive identity gives $De=g$. For compact smooth $h$,
\[
 q(g,h)=q(e,Dh)=\alpha_L\langle e,Dh\rangle_2
                   =\alpha_L\langle g,h\rangle_2.
\]
The form core extends this identity to $h=g$, contradicting the odd
minimum. Thus $c=0$, and injectivity of $C_L$ gives $[f]=0$.

Set $F=\widehat e$. It is even, entire of exponential type at most $L$,
and $F(0)\ne0$. The primitive identity gives
\[
 \widehat {K_Lf}(z)=
 \frac{\widehat f(z)-F(z)\ell(f)/\ell(e)}{z},\qquad z\ne0.
\]
At a zero of $F$, Fourier evaluation descends to a nonzero bounded
functional on the quotient. Its Riesz vector is an eigenvector with
eigenvalue $1/\overline z$, so selfadjointness makes $z$ real.
Conversely an eigenvector of the quotient with eigenvalue $\beta\ne0$
defines a distribution $\Lambda$ on $(-L,L)$ satisfying
$\Lambda(Dh)=\beta^{-1}\Lambda(h)$. Solving this constant-coefficient
first-order distributional equation gives
$\Lambda(h)=c\widehat h(\beta^{-1})$ with $c\ne0$. Since $\Lambda$
annihilates $e$, $F(\beta^{-1})=0$. This equation also shows the
eigenspace has dimension one. A double zero of $F$ would make evaluation
and its derivative descend simultaneously; differentiating the last
display gives a generalized eigenvector $w'$ satisfying
$(K_L-z^{-1})w'=-z^{-2}w$, where $w\ne0$ is the evaluation eigenvector.
Pairing with $w$ contradicts selfadjointness. All Fourier zeros are
therefore simple.

For the limiting comparison, take a compact smooth $h$ of integral one
and put $h_R(x)=R^{-1}h(x/R)$. Integration by parts gives
$q(h_R)\le R^{-2}\|h'\|_1^2\sum m_\gamma/\gamma^2$, while
$\|h_R\|_2^2=R^{-1}\|h\|_2^2$. Taking $R$ proportional to $L$ proves
$a_1(L)\to0$. Extend $M_L$ by zero to $H$. It is a positive contraction,
and for each fixed compact smooth $f$,
\[
 \|M_Lf\|_q^2\le q(M_Lf,f)=\alpha_L\|f\|_2^2\longrightarrow0.
\]
Density gives strong convergence $M_L\to0$. The increasing projections
$\Pi_L$ converge strongly to $I$, so Hilbert--Schmidt approximation gives
$C_L\to C$, $M_LC_L\to0$ and $C_LM_L\to0$ in
$\mathcal S_2$.

For $0\le x\le1-\kappa_L$,
\[
 (1-x)^{-1/2}-1
 =\frac{x}{\sqrt{1-x}(1+\sqrt{1-x})},\qquad
 |(1-x)^{1/2}-1|\le x.
\]
Apply these identities by functional calculus to the two blocks of
$B_L-C_L$. They give \eqref{eq:parity-error}, and the assumed positive
lower bound for $\kappa_L$ gives $B_L\to C$ in Hilbert--Schmidt norm
along the specified sequence.

Let $\beta_{L,j}>0$ be the positive eigenvalues of $B_L$ in decreasing
order. Min--max gives $\beta_{L,j}\to1/\gamma_j$ for every fixed $j$.
Reflection and Hilbert--Schmidt convergence also give
$\sum_j\beta_{L,j}^2\to\sum_j\gamma_j^{-2}$. Subtracting finitely
many convergent terms controls the remaining sums uniformly.
Hadamard factorization of the even exponential-type function $F$ gives
\[
 F(z)/F(0)=\prod_j(1-z^2\beta_{L,j}^2).
\]
For $|z|\le R$, the logarithms of sufficiently late factors have total
modulus at most $2R^2\sum_{j>J}\beta_{L,j}^2$. The uniform tail bound
and fixed-index convergence prove the asserted entire limit.
\end{proof}

\begin{proposition}[A subsequence criterion for simplicity]
\label{prop:liminf-rayleigh}
Let $K\in L^2(\mathbb R)$ be the inverse Fourier transform of $\Xi$,
and let $v_L$ be real even $L^2$-unit vectors in the local Weil form
domains with $v_L\to K/\|K\|_2$ in $L^2(\mathbb R)$.
The equality
\begin{equation}
 \liminf_{L\to\infty}\frac{q(v_L)}{a_2(L)}=0
 \label{eq:liminf-rayleigh}
\end{equation}
forces every nontrivial zeta zero to be simple. Equivalently, the
existence of a multiple zero forces a strictly positive lower limit
for this quotient. This applies to the original even prolate trial
of Connes, Consani and Moscovici.
\end{proposition}

\begin{proof}
Choose a sequence of windows along which
$r_L=q(v_L)/a_2(L)\to0$. Min--max gives $a_1\le q(v_L)<a_2$ eventually,
so the bottom line is simple. An odd bottom vector would be orthogonal
to the even trial and would give $q(v_L)\ge a_2$. Thus the bottom vector
$u_L$ is real and even. Choose its sign with positive overlap with $v_L$.
An eigenvector expansion gives
\[
 1-|\langle u_L,v_L\rangle_2|^2\le r_L,
 \qquad \|u_L-v_L\|_2^2\le2r_L.
\]
Hence $u_L\to K/\|K\|_2$. Furthermore
$a_e/a_o\le a_1/a_2\le r_L$, so the parity comparison applies along
this sequence and gives
$G_L=\widehat u_L/\widehat u_L(0)\to\Delta_\Gamma$ locally uniformly.
Plancherel gives $\widehat u_L\to\Xi/\|K\|_2$ in $L^2(\mathbb R)$.

Choose a small interval $I$ about zero on which $\Delta_\Gamma$ is
nonzero. With $c_L=\widehat u_L(0)$, the exact identity
\[
 c_L=\frac{\langle\widehat u_L,G_L\rangle_{L^2(I)}}
                {\|G_L\|_{L^2(I)}^2}
\]
shows that $c_L$ has a finite limit $c$. Thus
$c\Delta_\Gamma=\Xi/\|K\|_2$ almost everywhere on $I$, and continuity
and the identity theorem extend the equality to the whole plane.
The limit $c$ is nonzero because $\Xi(0)\ne0$, and evaluation at zero
gives $c=\Xi(0)/\|K\|_2$. Therefore
$\Delta_\Gamma=\Xi/\Xi(0)$. Since the factors of $\Delta_\Gamma$
use each ordinate once, every multiplicity is one.
\end{proof}

The estimate in the last proposition remains a question about the Weil
operator. It has the weaker subsequence quantifier displayed above;
neither positivity of the form nor the compact comparison proves it.
For the original prolate trial the available upper bound is
$q(v_{\log\lambda})\ll_\varepsilon
\lambda^\varepsilon(1-\nu_4)$, with
$1-\nu_4\asymp\lambda^9e^{-4\pi\lambda^2}$.
An upper bound for a trial does not provide the missing comparison with
$a_2$; the needed quantitative estimate is precisely the lower-limit
equality in the proposition.
